


\chapter{Introduction}

\begin{quotation}
\noindent\textbf{This guide describes an earlier version of the tool and is
kept as a historical record; it is no longer current.} It was written in
March 2020, when the script ran inside a Qt desktop window, could be
packaged into a standalone \texttt{dash-lineplot.exe} with PyInstaller, and
offered Plotly subplots and an x-axis slider. All of that has since been
removed: the Qt/PySide window and the packaged executable are gone (the
script serves its pages to the system browser directly), subplots were
removed outright, and the slider was replaced by the X/Y range boxes beside
each graph. JSON configuration files are also now accepted alongside Excel
workbooks, which this guide does not mention at all.
For the current, maintained guide, see \texttt{docs/userguide.md} in the
repository root -- where this guide and that one disagree, the Markdown
guide is correct. See \texttt{suggestedwork.md} for the tracked status of
bringing this LaTeX guide itself up to date.
\end{quotation}

The Python script \texttt{dash-lineplot.py} is a general plotting utility that aids in visualisation of captured or recorded data. It reads an Excel configuration file and forms a set of Dash data structures in \ac{HTML} pages. A Dash portal is created where the \ac{HTML} pages are served via a Flask server.
The user has full control of the graph sets that are rendered on the pages.
Figure~\ref{fig:examplePage} shows an example of the Dash portal with different graph sets being rendered on several tabs.


\begin{figure}[h]
\centering
\includegraphics[width=0.70\textwidth]{pic/examplePage}
\caption{Example pages rendered by the dash-lineplot.py utility Python script.
\label{fig:examplePage}}
\end{figure}
